% GENERATED FILE. Edit canonical lessons, then run npm run generate:books.
% canonical-source-hash: fnv1a64:9c713ef7d2625f45
% canonical-lessons: LA-C314-offendo, LA-C314-praeceps, LA-C314-praecipito, LA-C314-plaga, LA-C314-saucius, LA-R314-review-307-314

\chapter{Cāsus in viā}
\label{ch:la-casus-in-via}
\hlchaptermodality{314}

\begin{chapteropening}
\textbf{By the end of this chapter:} \emph{I can describe a fall or a knock: to strike against, headlong, to hurl down, a blow and wounded.}

Complete the last lesson of chapter 314: I can describe a fall or a knock: to strike against, headlong, to hurl down, a blow and wounded.
\end{chapteropening}

\section[offendō, offendere]{offendō, offendere — to strike against, to stumble; to offend}

\label{lesson:LA-C314-offendo}



\begin{quote}
Before the new one: say the Latin for to fall in, then the Latin for to fall down.
\end{quote}

\subsection*{You'll want to know: offendō, offendere}
\textbf{offendō, offendere} — \textquotedblleft{}to strike against, to stumble; to offend\textquotedblright{}.

\begin{quote}
Nocte pedem ad saxum offendī et cecidī.
\end{quote}

\begin{quote}
\emph{At night I struck my foot against a stone and fell.}
\end{quote}

\emph{offendō, offendere, offendī, offēnsum}. First a knock against something, then hurt feelings. English \emph{offend}.

\begin{cousinweb}
\textbf{offendō} joins \emph{ob-}, \textquotedblleft{}against\textquotedblright{}, and an old verb \emph{*fendō}, \textquotedblleft{}to strike\textquotedblright{} (EtymDB), also seen in \emph{dēfendō}, \textquotedblleft{}to ward off\textquotedblright{}. English \emph{offend} comes from it.
\end{cousinweb}

\begin{grammarlens}[title={What you've built}]
One more word for telling, explaining and describing at B1.
\end{grammarlens}

\subsection*{Your turn}
\begin{itemize}
\raggedright
  \item \emph{Say it:} \emph{offendō, offendere}
  \item \emph{Say it:} \emph{offendō, offendere}, once more
  \item \emph{Say it:} a sentence of your own with \emph{offendō, offendere}
  \item \emph{Recall:} say the Latin for to fall in, then the Latin for to fall down, then say \emph{offendō, offendere} again
\end{itemize}

\begin{tcolorbox}[breakable,colback=teal!4,colframe=teal!35!black,title={Before you move on}]
What does offendō, offendere mean? (\textquotedblleft{}To strike against, to stumble; to offend\textquotedblright{}.) Say it once more.
\end{tcolorbox}
\section[praeceps, praecipitis]{praeceps, praecipitis — headlong; steep}

\label{lesson:LA-C314-praeceps}



\begin{quote}
Before the new one: say the Latin for to fall down, then the Latin for to strike against.
\end{quote}

\subsection*{You'll want to know: praeceps, praecipitis}
\textbf{praeceps, praecipitis} — \textquotedblleft{}headlong; steep\textquotedblright{}.

\begin{quote}
Senex praeceps per scālās cecidit.
\end{quote}

\begin{quote}
\emph{The old man fell headlong down the stairs.}
\end{quote}

One ending, genitive \emph{praecipitis}: head first. Also 'steep': \emph{locus praeceps}.

\begin{cousinweb}
\textbf{praeceps} joins \emph{prae}, \textquotedblleft{}in front\textquotedblright{}, and \emph{caput}, \textquotedblleft{}head\textquotedblright{} (EtymDB): head first. English \emph{precipice} comes from its family.
\end{cousinweb}

\begin{grammarlens}[title={What you've built}]
One more word for telling, explaining and describing at B1.
\end{grammarlens}

\subsection*{Your turn}
\begin{itemize}
\raggedright
  \item \emph{Say it:} \emph{praeceps, praecipitis}
  \item \emph{Say it:} \emph{praeceps, praecipitis}, once more
  \item \emph{Say it:} a sentence of your own with \emph{praeceps, praecipitis}
  \item \emph{Recall:} say the Latin for to fall down, then the Latin for to strike against, then say \emph{praeceps, praecipitis} again
\end{itemize}

\begin{tcolorbox}[breakable,colback=teal!4,colframe=teal!35!black,title={Before you move on}]
What does praeceps, praecipitis mean? (\textquotedblleft{}Headlong; steep\textquotedblright{}.) Say it once more.
\end{tcolorbox}
\section[praecipitō, praecipitāre]{praecipitō, praecipitāre — to hurl headlong; to fall headlong}

\label{lesson:LA-C314-praecipito}



\begin{quote}
Before the new one: say the Latin for to strike against, then the Latin for headlong.
\end{quote}

\subsection*{You'll want to know: praecipitō, praecipitāre}
\textbf{praecipitō, praecipitāre} — \textquotedblleft{}to hurl headlong; to fall headlong\textquotedblright{}.

\begin{quote}
Equus subitō sē in fossam praecipitāvit.
\end{quote}

\begin{quote}
\emph{The horse suddenly threw itself headlong into the ditch.}
\end{quote}

\emph{praecipitō, praecipitāre, praecipitāvī, praecipitātum}, from \emph{praeceps}. English \emph{precipitate}.

\begin{cousinweb}
\textbf{praecipitō} is built on \emph{praeceps}, \textquotedblleft{}headlong\textquotedblright{} (EtymDB): to throw head first. English \emph{precipitate} comes from it.
\end{cousinweb}

\begin{grammarlens}[title={What you've built}]
One more word for telling, explaining and describing at B1.
\end{grammarlens}

\subsection*{Your turn}
\begin{itemize}
\raggedright
  \item \emph{Say it:} \emph{praecipitō, praecipitāre}
  \item \emph{Say it:} \emph{praecipitō, praecipitāre}, once more
  \item \emph{Say it:} a sentence of your own with \emph{praecipitō, praecipitāre}
  \item \emph{Recall:} say the Latin for to strike against, then the Latin for headlong, then say \emph{praecipitō, praecipitāre} again
\end{itemize}

\begin{tcolorbox}[breakable,colback=teal!4,colframe=teal!35!black,title={Before you move on}]
What does praecipitō, praecipitāre mean? (\textquotedblleft{}To hurl headlong; to fall headlong\textquotedblright{}.) Say it once more.
\end{tcolorbox}
\section[plāga, plāgae]{plāga, plāgae — a blow, a wound}

\label{lesson:LA-C314-plaga}



\begin{quote}
Before the new one: say the Latin for headlong, then the Latin for to hurl headlong.
\end{quote}

\subsection*{You'll want to know: plāga, plāgae}
\textbf{plāga, plāgae} — \textquotedblleft{}a blow, a wound\textquotedblright{}.

\begin{quote}
Plāga in capite gravis erat.
\end{quote}

\begin{quote}
\emph{The blow on the head was serious.}
\end{quote}

Feminine, with a long \emph{ā}. Do not confuse \emph{plaga}, 'a region' or 'a hunting net'. English \emph{plague} comes from it.

\begin{cousinweb}
\textbf{plāga} belongs with \emph{plangō}, \textquotedblleft{}to strike, to beat\textquotedblright{} (EtymDB), and matches Greek \emph{plēgē}, \textquotedblleft{}a blow\textquotedblright{} (Lewis \& Short). English \emph{plague}, a blow from heaven, comes from it.
\end{cousinweb}

\begin{grammarlens}[title={What you've built}]
One more word for telling, explaining and describing at B1.
\end{grammarlens}

\subsection*{Your turn}
\begin{itemize}
\raggedright
  \item \emph{Say it:} \emph{plāga, plāgae}
  \item \emph{Say it:} \emph{plāga, plāgae}, once more
  \item \emph{Say it:} a sentence of your own with \emph{plāga, plāgae}
  \item \emph{Recall:} say the Latin for headlong, then the Latin for to hurl headlong, then say \emph{plāga, plāgae} again
\end{itemize}

\begin{tcolorbox}[breakable,colback=teal!4,colframe=teal!35!black,title={Before you move on}]
What does plāga, plāgae mean? (\textquotedblleft{}A blow, a wound\textquotedblright{}.) Say it once more.
\end{tcolorbox}
\section[saucius, saucia, saucium]{saucius, saucia, saucium — wounded, hurt}

\label{lesson:LA-C314-saucius}



\begin{quote}
Before the new one: say the Latin for to hurl headlong, then the Latin for a blow.
\end{quote}

\subsection*{You'll want to know: saucius, saucia, saucium}
\textbf{saucius, saucia, saucium} — \textquotedblleft{}wounded, hurt\textquotedblright{}.

\begin{quote}
Medicus mīlitēs sauciōs cūrābat.
\end{quote}

\begin{quote}
\emph{The doctor was looking after the wounded soldiers.}
\end{quote}

Wounded but alive; \emph{vulnus} (taught) is the wound itself. Caesar counts the \emph{sauciī} after a battle.

\begin{cousinweb}
\textbf{saucius} is a Latin word whose origin is not settled (Lewis \& Short, EtymDB).
\end{cousinweb}

\begin{grammarlens}[title={What you've built}]
One more word for telling, explaining and describing at B1.
\end{grammarlens}

\subsection*{Your turn}
\begin{itemize}
\raggedright
  \item \emph{Say it:} \emph{saucius, saucia, saucium}
  \item \emph{Say it:} \emph{saucius, saucia, saucium}, once more
  \item \emph{Say it:} a sentence of your own with \emph{saucius, saucia, saucium}
  \item \emph{Recall:} say the Latin for to hurl headlong, then the Latin for a blow, then say \emph{saucius, saucia, saucium} again
\end{itemize}

\begin{tcolorbox}[breakable,colback=teal!4,colframe=teal!35!black,title={Before you move on}]
What does saucius, saucia, saucium mean? (\textquotedblleft{}Wounded, hurt\textquotedblright{}.) Say it once more.
\end{tcolorbox}
\section[(dialogue)]{Words from chapters 307–314 — a review}

\label{lesson:LA-R314-review-307-314}



\begin{quote}
Say the Latin for poison and a poisoner.
\end{quote}

\subsection*{Your turn}
Say these aloud:

\begin{itemize}
\raggedright
  \item poison, a poisoner, an ambush, a conspiracy, a traitor
  \item a judge, a trial, to judge, a law, law
  \item to accuse, a prosecutor, a supporter in court, a patron, evidence
  \item to swear, perjury, to lie, to prove, information
  \item to condemn, to acquit, guilty, innocent, fault
  \item to admit, to punish, punishment, exile, a chain
  \item to rush, a collapse, to collapse, to fall in, to fall down
  \item to strike against, headlong, to hurl headlong, a blow, wounded
\end{itemize}

\begin{tcolorbox}[breakable,colback=teal!4,colframe=teal!35!black,title={Before you move on}]
Poison? (\textbf{venēnum, venēnī}.) A poisoner? (\textbf{venēficus, venēficī}.) An ambush? (\textbf{īnsidiae, īnsidiārum}.) A conspiracy? (\textbf{coniūrātiō, coniūrātiōnis}.) A traitor? (\textbf{prōditor, prōditōris}.) A judge? (\textbf{iūdex, iūdicis}.) A trial? (\textbf{iūdicium, iūdiciī}.) To judge? (\textbf{iūdicō, iūdicāre}.) A law? (\textbf{lēx, lēgis}.) Law? (\textbf{iūs, iūris}.) To accuse? (\textbf{accūsō, accūsāre}.) A prosecutor? (\textbf{accūsātor, accūsātōris}.) A supporter in court? (\textbf{advocātus, advocātī}.) A patron? (\textbf{patrōnus, patrōnī}.) Evidence? (\textbf{testimōnium, testimōniī}.) To swear? (\textbf{iūrō, iūrāre}.) Perjury? (\textbf{periūrium, periūriī}.) To lie? (\textbf{mentior, mentīrī}.) To prove? (\textbf{arguō, arguere}.) Information? (\textbf{indicium, indiciī}.) To condemn? (\textbf{damnō, damnāre}.) To acquit? (\textbf{absolvō, absolvere}.) Guilty? (\textbf{nocēns, nocentis}.) Innocent? (\textbf{innocēns, innocentis}.) Fault? (\textbf{culpa, culpae}.) To admit? (\textbf{fateor, fatērī}.) To punish? (\textbf{pūniō, pūnīre}.) Punishment? (\textbf{supplicium, suppliciī}.) Exile? (\textbf{exsilium, exsiliī}.) A chain? (\textbf{vinculum, vinculī}.) To rush? (\textbf{ruō, ruere}.) A collapse? (\textbf{ruīna, ruīnae}.) To collapse? (\textbf{corruō, corruere}.) To fall in? (\textbf{collābor, collābī}.) To fall down? (\textbf{dēcidō, dēcidere}.) To strike against? (\textbf{offendō, offendere}.) Headlong? (\textbf{praeceps, praecipitis}.) To hurl headlong? (\textbf{praecipitō, praecipitāre}.) A blow? (\textbf{plāga, plāgae}.) Wounded? (\textbf{saucius, saucia, saucium}.)
\end{tcolorbox}
